\clearpage
\section{Conclusions}

The final project preserves the \textbf{alarm semantics} of assignment 3 and adds
the distributed mechanisms required by assignment 4: \textbf{cluster startup},
\textbf{remote discovery}, \textbf{singleton state ownership}, and \textbf{safe
recovery after recreation}.

The resulting system runs on multiple cluster nodes without changing the alarm
rules of the non-clustered version. Sensors, keypad, and control logic may run
on different nodes, but communication remains message-based and the alarm state
is still owned by one logical control entity.

The main architectural trade-off is explicit: the project does not introduce
persistence or replicated alarm state. If the control unit is recreated, the
previous in-memory state is lost and the system returns to
\texttt{STARTUP\_RECOVERY}. This behavior matches the specification and keeps the
design aligned with the actual assignment scope.
